\documentclass[10pt,a4paper]{amsart}
\usepackage[margin=19mm]{geometry}
\usepackage{amsmath,amssymb,mathtools}
\usepackage[scr=rsfs,cal=euler]{mathalfa}
\usepackage{xfrac}
\usepackage[unicode,hidelinks,bookmarks=false]{hyperref}
\newcommand{\ie}{i.e.\ }
\newcommand{\cf}{cf.\ }
\newcommand{\resp}{respectively\ }
\newcommand{\Subsection}{\subsection}
\providecommand{\nobreakdash}{\nobreak\hbox{-}}
\setlength{\emergencystretch}{2em}
\multlinegap0pt
\usepackage{bm}
\usepackage{bbm}
\usepackage{dsfont}
\usepackage{xspace}
\usepackage{cancel}
\usepackage{mathtools}
\usepackage{amsthm}
\makeatletter
\@ifundefined{theorem}{\newtheorem{theorem}{Theorem}}{}
\@ifundefined{corollary}{\newtheorem{corollary}[theorem]{Corollary}}{}
\@ifundefined{lemma}{\newtheorem{lemma}[theorem]{Lemma}}{}
\@ifundefined{proposition}{\newtheorem{proposition}[theorem]{Proposition}}{}
\makeatother
\newcommand\E{{\mathbb E}}
\newcommand\Pbb{{\mathbb{P}}}
\newcommand\U{{\mathrm U}}
\newcommand\Z{{\mathbb{Z}}}
\title[The heat-kernel master field on $\mathbb{Z}^d$ at strong cou]{The heat-kernel master field on $\mathbb{Z}^d$ at strong coupling}
\author{Thibaut Lemoine}
\date{Source: arXiv:2606.28945v2 (CC BY 4.0)}
\begin{document}
\maketitle

\noindent\emph{Source and licence.} The heat-kernel master field on $\mathbb{Z}^d$ at strong coupling. Version arXiv:2606.28945v2 (CC BY 4.0), distributed under the Creative Commons Attribution 4.0 International licence, as recorded in the arXiv licence metadata for this exact version and shown on the abstract page https://arxiv.org/abs/2606.28945v2. Full rights notice: licenses/arxiv-2606.28945-CC-BY-4.0.txt.\par\smallskip

\noindent\textit{Editorial scope.} Counted unit: the Theorem whose source label is \texttt{thm:stable-hk-block} in the article (source0.tex lines 1489--1503, source label \texttt{thm:stable-hk-block}), delivered together with the article's own complete original proof of it (source0.tex lines 1505--1572, one \texttt{proof} environment of 68 lines). No mathematics is written, repaired or re-derived: statement and proof are the article's own text line for line. Required context retained uncounted: eq:spectral-decay (lines 358--360); eq:split-selection (lines 885--889); lem:central-cutoff (lines 1264--1270); thm:pinned-block (lines 1377--1386); lem:minimal-support-size (lines 1392--1398); prop:nonstable-polymer-tail (lines 1440--1460); cor:coarse-nonspectral-block (lines 2039--2047). Every definition, display and label of those lines is retained unchanged. Omitted from this package and not redistributed: 1--357, 361--884, 890--1263, 1271--1376, 1387--1391, 1399--1439, 1461--1488, 1504--1504, 1573--2038, 2048--3501. Nothing inside a retained excerpt is omitted or altered: every retained body is delivered exactly as the article's lines are written, so no omission receipt of any kind accompanies this package. Citations: the retained excerpts carry no citation command at all, so no bibliography entry of the article is transcribed and no reference is redistributed. Reference adaptation: none was needed and none was made; every reference inside the delivered unit resolves inside it, so no unresolved reference or citation is produced. Printed slips kept literal: no printed word, symbol, hypothesis or number of the counted statement or of its proof was altered. Layout and class: the article's own class is \texttt{amsart} with packages including geometry, microtype, amsmath, amsfonts, amssymb, amsthm, mathrsfs, graphicx, hyperref, enumitem; the wrapper is the lane's pinned \texttt{amsart} editorial wrapper at 10pt on A4, which restates the theorem-like environments the retained text needs and adds the article's own macro definitions that the retained text needs (\texttt{\textbackslash E}, \texttt{\textbackslash Pbb}, \texttt{\textbackslash U}, \texttt{\textbackslash Z}), with extra wrapper packages (bm, bbm, dsfont, xspace, cancel, mathtools). The delivered numbering is the unit's own. Assets: no retained span contains a figure environment, a graphics-inclusion command, a PDF-page inclusion, a table or a TikZ picture; no image file is copied into this package, the package declares no asset dependency and no image file is redistributed with it. Licence: version arXiv:2606.28945v2 of this article is distributed under the Creative Commons Attribution 4.0 International licence, as recorded in the arXiv licence metadata for this exact version and shown on the abstract page https://arxiv.org/abs/2606.28945v2. A full rights notice is delivered at licenses/arxiv-2606.28945-CC-BY-4.0.txt. Characters: every retained excerpt is pure ASCII, so the delivered engine, XeLaTeX with its default Latin Modern text font, reports no missing character.
\begin{equation}\label{eq:spectral-decay}
\exp\left(-\frac{T}{2N}C_2(\lambda)\right)\leq \exp\left(-\frac{T}{4}\mathcal E(\lambda)\right)\exp\left(\frac{T m(\lambda)^2}{2N^2}\right),
\end{equation}

\begin{equation}\label{eq:split-selection}
\partial^*q=0,
\qquad
\partial^*s+j_L=0.
\end{equation}

\begin{lemma}\label{lem:central-cutoff}
Fix $0<\varepsilon_d<(4d)^{-1}$.  Let $Y$ be a connected active support and let $M=M(\alpha_Y)$ be the total projective mass of the reduced plaquette decoration on $P(Y)$.  For every fixed loop family $L$ there exists $N_0(L,d)<\infty$ such that, whenever $N\ge N_0(L,d)$ and $M\leq \varepsilon_d N$, the central selection rule \eqref{eq:split-selection} holds on the part of the stable expansion supported on $Y$. Moreover, suppose that the corresponding non-spectral stable contribution on $Y$, after summing over all internal local-channel variables of total projective mass $M$ but before inserting the reduced heat-kernel factor, is bounded by $C_LC^{|Y|+M}$.  Then the complementary stable range $M>\varepsilon_dN$ contributes at most
\begin{equation}\label{eq:stable-large-mass-tail}
C_L C^{|Y|}\exp\{-(cT-\log C)\varepsilon_d N\}
\end{equation}
to the corresponding support, for a constant $c>0$ independent of $N$ and $\Lambda$.
\end{lemma}

\begin{theorem}\label{thm:pinned-block}
There exists $C_d<\infty$, depending only on $d$, such that the following holds.  Let $\Lambda\subset\Z^d$ be finite, let $L$ be fixed, and let $Y$ be a connected pinned support.  For every $N$ and every $M\le N/2$,
\begin{equation}
\label{eq:pinned-block-bound}
\left\|\mathcal O_{Y,N}^{\mathrm{red},L}\right\|_{M,1}
\le
C_L C_d^{|Y|+M}N^{-\delta(Y,L)},
\end{equation}
where $\delta(Y,L)\ge0$, and where $\delta(Y,L)\ge1$ if $b(Y,L)\ge2$.  Moreover, for fixed $Y$ and $M$, the signed block coefficient $\mathcal O_{Y,N}^{\mathrm{red},L}[M]$ admits a controlled asymptotic expansion to all orders in powers of $1/N$, uniformly in the ambient finite volume.
\end{theorem}

\begin{lemma}
\label{lem:minimal-support-size}
There exist $a_d>0$ and $C_L<\infty$ such that, for every connected pinned support $Y$ and every plaquette decoration $\lambda_Y:P(Y)\to\widehat{\U(N)}$,
\[
\mathcal{E}(\lambda_Y)\ge a_d |Y|-C_L .
\]
\end{lemma}

\begin{proposition}\label{prop:nonstable-polymer-tail}
Fix $A\ge1$ and assume $T>4\log A$.  For every connected pinned support $Y$ and every plaquette label field $\alpha=(\alpha_p)_{p\in P(Y)}$, let $B_N^L(Y;\alpha)$ be a non-negative quantity satisfying
\begin{equation}
\label{eq:nonspectral-growth-assumption}
B_N^L(Y;\alpha)
\le
C_LC^{|Y|}A^{M(\alpha)},
\qquad
M(\alpha)=\sum_{p\in P(Y)}m(\alpha_p),
\end{equation}
with constants independent of $N$ and of the finite volume.  Then, for some $C_T<\infty$ and $\gamma(A,T)>0$,
\begin{equation}\label{eq:nonstable-polymer-tail-under-P}
\widehat\E_{Y,T,N}\left[
B_N^L(Y;\alpha)
\mathbf 1_{\{M(\alpha)>N/2\}}
\right]
\le
C_L C_T^{|Y|}e^{-\gamma(A,T)N},
\end{equation}
whenever $T>4\log A$, uniformly in $\Lambda$ and $N$.
\end{proposition}

\begin{corollary}[Coarse non-spectral block bound]\label{cor:coarse-nonspectral-block}
There are constants $C_d<\infty$ and $C_L<\infty$ such that, for every connected pinned support $Y$ and every reduced plaquette-label field $\alpha_Y$ of arbitrary total projective mass
$M=M(\alpha_Y)$,
\begin{equation}\label{eq:coarse-nonspectral-block}
\left|\mathcal O_{Y,N}^{\mathrm{red},L}(\alpha_Y)\right|
\le C_LC_d^{|Y|+M}.
\end{equation}
The same estimate holds after arbitrary determinant shifts are reinserted, provided the central charge selection rule is imposed.
\end{corollary}

\begin{theorem}
\label{thm:stable-hk-block}
There exist constants $C,c>0$, depending only on $d$, with the following property.  For every fixed loop family $L$, there is $C_L<\infty$ such that, for every connected pinned support $Y$,
\begin{equation}
\label{eq:stable-hk-block}
\widehat\E_{Y,T,N}\left[
\mathcal O_{Y,N}^{\mathrm{abs},\mathrm{st},L}(\lambda_Y)
\right]
\le
C_L\exp\{-\mu(T)|Y|\},
\qquad
\mu(T)=cT-\log C .
\end{equation}

\end{theorem}

\begin{proof}
We estimate the expectation in \eqref{eq:stable-hk-block} by using the determinant-reduced block observable and the product reference law $\widehat\Pbb_{Y,T,N}$. We first restrict to the stable small-mass range $M=\sum_{p\in P(Y)}m(\alpha_p)\le \varepsilon_d N,$ where $\varepsilon_d$ is as in Lemma~\ref{lem:central-cutoff}.  On this range the central selection rule splits as
\[
\partial^*q=0,
\qquad
\partial^*s+j_L=0.
\]
Thus the determinant characters cancel from the local-channel coefficient; the determinant variables remain only in the reference law $\widehat\Pbb_{T,N}$.

We now estimate the contribution at fixed projective mass $M$, after the local-channel variables have been resummed and the reduced labels are summed.  Using the heat-kernel domination bound~\eqref{eq:spectral-decay} plaquettewise, and absorbing the factor $\exp\{\frac{T}{2}\sum_p m(\alpha_p)^2/N^2\}$ into the constants on the range $M\le\varepsilon_dN$, there is a constant $c>0$ such that, after changing constants,
\[
\prod_{p\in P(Y)}
\exp\left\{-\frac{T}{2N}C_2(q_p\mathbf 1_N+\alpha_p)\right\}
\le C_L\exp\{-cT(M+Q)\}.
\]
Here and below,
\[
Q:=\sum_{p\in P(Y)}q_p^2.
\]
The possible factor coming from the cross term $2q_ps_p/N$ is absorbed by decreasing $c$, using $M\le\varepsilon_dN$ and the elementary inequality
\[
\frac{2|q_p|m_p}{N}
\le \frac{q_p^2}{2}+\frac{2m_p^2}{N^2}.
\]

For fixed $M$ and $q_Y$, Theorem~\ref{thm:pinned-block} already includes the sum over all reduced label fields of total mass $M$, with all local channel variables resummed at each fixed label field:
\[
\left\|\mathcal O_{Y,N}^{\mathrm{red},L}\right\|_{M,1}
\le C_LC_d^{|Y|+M}N^{-\delta(Y,L)} .
\]
Discarding the harmless factor $N^{-\delta(Y,L)}\le1$, the absolute contribution of all decorations with reduced mass $M$ and fixed determinant shifts $q_Y$ is therefore bounded by
\[
C_L C^{|Y|+M}\exp\{-cT(M+Q)\} .
\]
For every plaquette decoration, Lemma~\ref{lem:minimal-support-size} gives
\[
\mathcal E(\lambda_Y)=M+Q\ge a_d|Y|-C_L.
\]
Hence, after changing constants,
\[
C_L C^{|Y|+M}e^{-cT(M+Q)}
\le
C_L e^{-cTa_d|Y|/2}
C^{|Y|+M}e^{-cT(M+Q)/2}.
\]
The remaining sum over $M$ and over determinant shifts $q_Y\in\Z^{P(Y)}$ is bounded by $C^{|Y|}$ for $T$ large enough: the reduced labels have already been summed in $\left\|\mathcal O_{Y,N}^{\mathrm{red},L}\right\|_{M,1}$, the residual factor $C^M e^{-cTM/2}$ is summable, and the determinant shifts are controlled by the product Gaussian sum $\sum_q e^{-cTq^2/2}<\infty$.  Thus
\[
\widehat\E_{Y,T,N}\left[
\mathcal O_{Y,N}^{\mathrm{abs},\mathrm{st},L}(\lambda_Y)
\mathbf 1_{\{M(\alpha_Y)\le\varepsilon_dN\}}
\right]
\le
C_L\exp\{-(c'T-\log C')|Y|\}.
\]

It remains to treat the complement of the stable small-mass range.  We split it into
\[
\varepsilon_dN<M\le N/2
\qquad\text{and}\qquad
M>N/2.
\]
In the first range, Theorem~\ref{thm:pinned-block} gives the same non-spectral majorant $C_LC^{|Y|+M}$, and Lemma~\ref{lem:central-cutoff} gives an exponentially small cutoff in $N$ after insertion of the reduced heat-kernel factor.  The determinant variables are again controlled by Gaussian tails, and Lemma~\ref{lem:minimal-support-size} extracts the same support decay from $\mathcal E(\lambda_Y)$.  In the second range, Corollary~\ref{cor:coarse-nonspectral-block} gives the pointwise non-spectral bound
$C_LC^{|Y|}A^{M}$ for a fixed constant $A=A(d)$; the tilted tail estimate of Proposition~\ref{prop:nonstable-polymer-tail} therefore applies directly.  For $T$ large enough both complementary contributions are bounded by
\[
C_L\exp\{-(c''T-\log C'')|Y|\}
\]
after changing the constants, and are in fact exponentially small in $N$.  Combining this with the small-mass estimate proves~\eqref{eq:stable-hk-block}.
\end{proof}

\end{document}
